% GENERATED FILE. Edit canonical lessons, then run npm run generate:books.
% canonical-source-hash: fnv1a64:4b5d5f6754141f0c
% canonical-lessons: LA-C547-imago, LA-C547-effigies, LA-C547-figura, LA-C547-tessella, LA-C547-caelo

\chapter{Images and art}
\label{ch:la-images-and-art}
\hlchaptermodality{547}

\begin{chapteropening}
\textbf{By the end of this chapter:} \emph{I can talk about portraits, shapes and mosaics.}

Complete the last lesson of chapter 547: I can talk about portraits, shapes and mosaics.
\end{chapteropening}

\section[imāgō, imāginis]{imāgō, imāginis — an image, likeness, portrait}

\label{lesson:LA-C547-imago}



\begin{quote}
Before the new one: say the Latin for a prologue, then the Latin for a mime dancer.
\end{quote}

\subsection*{You'll want to know: imāgō, imāginis}
\textbf{imāgō, imāginis} — \textquotedblleft{}an image, likeness, portrait\textquotedblright{}.

\begin{quote}
In ātriō imāginēs maiōrum sunt.
\end{quote}

\begin{quote}
\emph{In the hall there are portraits of the ancestors.}
\end{quote}

Feminine, third declension. Also the wax masks of ancestors kept in the hall. English \emph{image}.

\begin{cousinweb}
\textbf{imāgō} is usually linked with \emph{imitor} \textquotedblleft{}to copy\textquotedblright{} (Lewis \& Short; EtymDB). English \emph{image} and \emph{imagine} come from it.
\end{cousinweb}

\begin{grammarlens}[title={What you've built}]
One more word for telling, explaining and describing at B1.
\end{grammarlens}

\subsection*{Your turn}
\begin{itemize}
\raggedright
  \item \emph{Say it:} \emph{imāgō, imāginis}
  \item \emph{Say it:} \emph{imāgō, imāginis}, once more
  \item \emph{Say it:} a sentence of your own with \emph{imāgō, imāginis}
  \item \emph{Recall:} say the Latin for a prologue, then the Latin for a mime dancer, then say \emph{imāgō, imāginis} again
\end{itemize}

\begin{tcolorbox}[breakable,colback=teal!4,colframe=teal!35!black,title={Before you move on}]
What does imāgō, imāginis mean? (\textquotedblleft{}An image, likeness, portrait\textquotedblright{}.) Say it once more.
\end{tcolorbox}
\section[effigiēs, effigiēī]{effigiēs, effigiēī — a likeness, effigy}

\label{lesson:LA-C547-effigies}



\begin{quote}
Before the new one: say the Latin for a mime dancer, then the Latin for an image.
\end{quote}

\subsection*{You'll want to know: effigiēs, effigiēī}
\textbf{effigiēs, effigiēī} — \textquotedblleft{}a likeness, effigy\textquotedblright{}.

\begin{quote}
Effigiēs avī in ātriō stat.
\end{quote}

\begin{quote}
\emph{A likeness of grandfather stands in the hall.}
\end{quote}

Feminine, fifth declension, from \emph{effingō} — to shape, portray. English \emph{effigy}.

\begin{cousinweb}
\textbf{effigiēs} is built on \emph{effingō} \textquotedblleft{}to form, portray\textquotedblright{}, from \emph{ex-} and \emph{fingō} \textquotedblleft{}to shape\textquotedblright{} (EtymDB; Lewis \& Short). English \emph{effigy} comes from it.
\end{cousinweb}

\begin{grammarlens}[title={What you've built}]
One more word for telling, explaining and describing at B1.
\end{grammarlens}

\subsection*{Your turn}
\begin{itemize}
\raggedright
  \item \emph{Say it:} \emph{effigiēs, effigiēī}
  \item \emph{Say it:} \emph{effigiēs, effigiēī}, once more
  \item \emph{Say it:} a sentence of your own with \emph{effigiēs, effigiēī}
  \item \emph{Recall:} say the Latin for a mime dancer, then the Latin for an image, then say \emph{effigiēs, effigiēī} again
\end{itemize}

\begin{tcolorbox}[breakable,colback=teal!4,colframe=teal!35!black,title={Before you move on}]
What does effigiēs, effigiēī mean? (\textquotedblleft{}A likeness, effigy\textquotedblright{}.) Say it once more.
\end{tcolorbox}
\section[figūra, figūrae]{figūra, figūrae — a shape, form, figure}

\label{lesson:LA-C547-figura}



\begin{quote}
Before the new one: say the Latin for an image, then the Latin for a likeness.
\end{quote}

\subsection*{You'll want to know: figūra, figūrae}
\textbf{figūra, figūrae} — \textquotedblleft{}a shape, form, figure\textquotedblright{}.

\begin{quote}
Pictor figūram equī in mūrō pingit.
\end{quote}

\begin{quote}
\emph{The painter paints the shape of a horse on the wall.}
\end{quote}

Feminine, first declension, from \emph{fingō}. English \emph{figure}.

\begin{cousinweb}
\textbf{figūra} is built on \emph{fingō} \textquotedblleft{}to shape, mould\textquotedblright{} (Lewis \& Short); EtymDB traces \emph{fingō} to a Proto-Indo-European root meaning \textquotedblleft{}to knead, shape\textquotedblright{}. English \emph{figure} comes from it.
\end{cousinweb}

\begin{grammarlens}[title={What you've built}]
One more word for telling, explaining and describing at B1.
\end{grammarlens}

\subsection*{Your turn}
\begin{itemize}
\raggedright
  \item \emph{Say it:} \emph{figūra, figūrae}
  \item \emph{Say it:} \emph{figūra, figūrae}, once more
  \item \emph{Say it:} a sentence of your own with \emph{figūra, figūrae}
  \item \emph{Recall:} say the Latin for an image, then the Latin for a likeness, then say \emph{figūra, figūrae} again
\end{itemize}

\begin{tcolorbox}[breakable,colback=teal!4,colframe=teal!35!black,title={Before you move on}]
What does figūra, figūrae mean? (\textquotedblleft{}A shape, form, figure\textquotedblright{}.) Say it once more.
\end{tcolorbox}
\section[tessella, tessellae]{tessella, tessellae — a small cube (of mosaic)}

\label{lesson:LA-C547-tessella}



\begin{quote}
Before the new one: say the Latin for a likeness, then the Latin for a shape.
\end{quote}

\subsection*{You'll want to know: tessella, tessellae}
\textbf{tessella, tessellae} — \textquotedblleft{}a small cube (of mosaic)\textquotedblright{}.

\begin{quote}
Faber pavīmentum ex multīs tessellīs facit.
\end{quote}

\begin{quote}
\emph{The craftsman makes the floor out of many little cubes.}
\end{quote}

Feminine, a diminutive of \emph{tessera} — a cube, a die. English \emph{tessellated}.

\begin{cousinweb}
\textbf{tessella} is a diminutive of \emph{tessera} \textquotedblleft{}a cube, die\textquotedblright{} (Lewis \& Short). English \emph{tessellated} comes from it.
\end{cousinweb}

\begin{grammarlens}[title={What you've built}]
One more word for telling, explaining and describing at B1.
\end{grammarlens}

\subsection*{Your turn}
\begin{itemize}
\raggedright
  \item \emph{Say it:} \emph{tessella, tessellae}
  \item \emph{Say it:} \emph{tessella, tessellae}, once more
  \item \emph{Say it:} a sentence of your own with \emph{tessella, tessellae}
  \item \emph{Recall:} say the Latin for a likeness, then the Latin for a shape, then say \emph{tessella, tessellae} again
\end{itemize}

\begin{tcolorbox}[breakable,colback=teal!4,colframe=teal!35!black,title={Before you move on}]
What does tessella, tessellae mean? (\textquotedblleft{}A small cube (of mosaic)\textquotedblright{}.) Say it once more.
\end{tcolorbox}
\section[caelō, caelāre]{caelō, caelāre — to engrave, carve in relief}

\label{lesson:LA-C547-caelo}



\begin{quote}
Before the new one: say the Latin for a shape, then the Latin for a small cube.
\end{quote}

\subsection*{You'll want to know: caelō, caelāre}
\textbf{caelō, caelāre} — \textquotedblleft{}to engrave, carve in relief\textquotedblright{}.

\begin{quote}
Faber pōculum argenteum caelat.
\end{quote}

\begin{quote}
\emph{The smith engraves a silver cup.}
\end{quote}

\emph{caelō, caelāre, caelāvī, caelātum}: raised work on silver, gold or ivory. Not to be confused with \emph{cēlō} — to hide.

\begin{cousinweb}
\textbf{caelō} is built on \emph{caelum} \textquotedblleft{}an engraver’s chisel\textquotedblright{}, a different word from \emph{caelum} \textquotedblleft{}sky\textquotedblright{} (Lewis \& Short).
\end{cousinweb}

\begin{grammarlens}[title={What you've built}]
One more word for telling, explaining and describing at B1.
\end{grammarlens}

\subsection*{Your turn}
\begin{itemize}
\raggedright
  \item \emph{Say it:} \emph{caelō, caelāre}
  \item \emph{Say it:} \emph{caelō, caelāre}, once more
  \item \emph{Say it:} a sentence of your own with \emph{caelō, caelāre}
  \item \emph{Recall:} say the Latin for a shape, then the Latin for a small cube, then say \emph{caelō, caelāre} again
\end{itemize}

\begin{tcolorbox}[breakable,colback=teal!4,colframe=teal!35!black,title={Before you move on}]
What does caelō, caelāre mean? (\textquotedblleft{}To engrave, carve in relief\textquotedblright{}.) Say it once more.
\end{tcolorbox}
